% Cover letter for the (new) IP&M submission of the Sieve manuscript.
% All factual statements are drawn from paper-journal-v2 (2026-10-01,
% commit e870b0d, 47 pp) and its auto-generated numbers.tex macros.
\documentclass[12pt]{article}
\usepackage[T1]{fontenc}
\usepackage[margin=1in]{geometry}
\usepackage{microtype}
\usepackage{parskip}
\pagestyle{empty}
\setlength{\parskip}{0.55em}

\newcommand{\PaperTitle}{Sieve: Budgeted Question-Conditioned
Information Selection for Training-Free Long-Context Compression}

\begin{document}

\begin{flushleft}
October 2, 2026

\bigskip
The Editors\\
Information Processing \& Management
\end{flushleft}

\bigskip
\noindent\textbf{Re: Submission of a Full Length Article --- ``\PaperTitle{}''}

\bigskip
Dear Editors,

Please consider the manuscript referenced above for publication in
\emph{Information Processing \& Management} as a \textbf{Full Length
Article}. I am the sole author of the manuscript and serve as its
corresponding author.

I would like to be transparent about this manuscript's history. An
earlier version was previously submitted to \emph{IP\&M} and was
declined, with the editorial assessment that the results appeared
premature. The present manuscript is a substantially reframed and
extended version, prepared specifically to address that assessment:
the contribution is now stated as a problem formulation rather than a
single system, the evidence base has been scaled and deepened with
three new experiment families, and the two protocols that previously
rested on the thinnest samples have been re-run at size. A summary of
what changed follows; I am submitting it as a new submission and leave
the editorial decision on how to handle the history entirely to you.

The manuscript formulates long-context compression as \emph{budgeted
question-conditioned information selection}: given a fixed word
budget, choose the subset of sentences that jointly maximizes
question-relevance coverage, positional reliability, and diversity,
with the selector's own computational cost treated as a first-class
constraint rather than an afterthought. This formulation, to our
knowledge absent from prior work even where its ingredients appear as
isolated heuristics, is stated formally, and a twelve-selector design
space table locates every prior system along seven dimensions
(question conditioning, positional awareness, redundancy control,
model-free inference, training-free operation, sentence-level intact
units, explicit budget with measured cost). The unoccupied corner the
table exposes is the one this paper measures: a selector that is
simultaneously question-conditioned, position-aware,
redundancy-controlling, and entirely model-free.

The vehicle of the study is \emph{Sieve}, the instantiation of that
formulation in this corner: IDF-weighted coverage of query content
terms, a U-shaped positional prior, and a redundancy penalty applied
greedily under the explicit budget, with selected sentences emitted in
their original order. Its preprocessing is a linear-time pass over the
context, its budgeted selection is $O(k\,n\,\bar{s})$, and the full
selector runs in tens of milliseconds on one CPU core (50.9--54.3\,ms
at a 6{,}000-word context), with no model parameters, no training,
and no dependency beyond NumPy.

The evaluation is organized as a three-way quality--compression--cost
frontier study, and it is the part that has grown the most since the
earlier version. First, the reader-free evidence-recall protocol now
spans three evidence-annotated benchmarks --- QASPER, HotpotQA
(distractor setting), and 2WikiMultihopQA --- 1{,}600 records in
total: \emph{Sieve} retains more gold evidence than every
question-agnostic baseline on all three datasets (margins over head
truncation of 12.8, 44.8, and 22.8 recall points at $4\times$), while
BM25 sentence ranking, the strongest question-conditioned model-free
comparator, exceeds the full selector by 6.8 points on QASPER --- a
result the manuscript reports as measured and reads through its own
ablations as a structural trade of raw recall for positional
reliability. Second, the end-task protocol is new since the earlier
submission and is now run at size: seven selection methods
(full-context, head, random, stride, TextRank, BM25, \emph{Sieve}) at
$n{=}100$ each under a production LLM reader at the $4\times$
operating point, with paired bootstrap intervals throughout. The
selector-level ordering reappears end to end: the BM25 reference
retains 76.0\% of full-context F1 against 61.8\% for \emph{Sieve}
(paired gap 4.8 points, interval excluding zero), while \emph{Sieve}
remains statistically matched with truncation and the question-blind
floor (paired gap +4.0, interval crossing zero) --- the same
trade-off, now consistent at both the selection and the answer layer.
Third, a full configuration--budget--dataset ablation matrix isolates
each component's contribution: the relevance term alone is a
near-oracle for evidence recovery (43.6\% versus 36.2\% for the full
configuration on QASPER), the positional prior's price is a
dataset-stable six-to-seven paired points, and all five selector
configurations run in the same milliseconds band, so the trade is
paid in information units rather than in compute. Fourth, the
reader-side positional profile that motivates the U-shaped prior is
no longer borrowed from cited work: it is re-measured under our own
reader at $n{=}50$ paired records across five evidence positions, and
reported flat at this context scale (median ${\sim}1{,}000$ words),
which the manuscript reports as a finding about scope: the prior
remains a selector-level design choice whose cost the ablation
prices. Fifth, a GPT-2 perplexity gate, measured on the same hardware
as every other method, reaches 28.7\% evidence recall at
255--320$\times$ \emph{Sieve}'s measured cost; and the correlation
between evidence recall and end-task answer F1 is re-estimated at the
new sample size (pooled Pearson 0.22 over 600 pairs), so the two
protocols are reported separately throughout.

I believe the manuscript fits the scope of \emph{Information Processing
\& Management}. It is a quantitative study of question-conditioned
information selection: it builds directly on classical
information-retrieval machinery (IDF weighting, BM25, maximal marginal
relevance), evaluates selection quality with recall protocols over
annotated evidence, connects selection decisions to the processing and
cost of managing long contexts in information systems, and reports
quality, compression, and compressor cost for every method so that
practitioners have a measured basis for deciding when statistical
selection suffices and when a hosted scorer earns its infrastructure.

The manuscript substantially extends an earlier conference-format
study of the same method, as described in the Introduction: a
three-dataset evaluation, the BM25 baseline, the same-hardware
comparison against a neural perplexity gate, the formal layer
(formulation, propositions, cost statements), the correlation study,
the end-to-end seven-method protocol at $n{=}100$, the
configuration--budget--dataset ablation matrix, and the reader-side
positional profile.

I confirm that: the manuscript is original and has not been published
previously in any form; it is not under consideration by any other
journal; I have no known competing financial interests or personal
relationships that could have appeared to influence the work reported;
and the work received no specific funding. The benchmarks used are
public (QASPER, HotpotQA, 2WikiMultihopQA), and the released
repository reproduces every table and figure from raw checkpoints.

Thank you for considering the manuscript. I look forward to the
reviewers' comments.

\bigskip
Sincerely,

\bigskip
Chang Tan\\
Undergraduate Student, School of Economics and Management\\
Dalian University of Technology\\
No.\,2 Linggong Road, Ganjingzi District, Dalian, Liaoning 116024,
China\\
E-mail: aquars43@foxmail.com

\end{document}
